\documentclass[11pt]{article}

\usepackage[margin=1in]{geometry}
\usepackage{amsmath,amssymb}
\usepackage{booktabs}
\usepackage{array}
\usepackage{longtable}
\usepackage{hyperref}

\title{GitHub-State Comparison and Continuum-Foundations Synchronization Release for HAOS-IIP}
\author{Tomislav Rupi\'c \\ Independent Researcher}
\date{April 2026}

\begin{document}
\maketitle

\begin{abstract}
This paper freezes the repository-comparison state of HAOS-IIP immediately before the next public synchronization push. The live GitHub branch remains at commit \texttt{cf5d48f}, which is the frozen geometry-emergence release \texttt{47.2}. The local repository, however, now contains a coherent continuum-facing tranche that has not yet been published: strengthened scalar operator-convergence controls, harmonic / coexact vector-sector separation work, the bounded HAOS-to-harmonic theorem ladder from \texttt{T1} through \texttt{T8}, the first active-sector transport machinery for the derived coexact physical branch, and two already prepared synthesis papers (\texttt{48.1} and \texttt{49.1}). This paper does not claim new physics beyond those artifacts. Its role is narrower and operationally important. It records exactly what is live on GitHub, exactly what exists locally above that live state, what remains open, what belongs in the next synchronized release, and what should still be excluded because it is unfinished or unrelated. The result is release \texttt{50.1}: a bounded synchronization paper that closes the comparison gap between public GitHub state and local repository state without overstating unfinished disorder-threshold work.
\end{abstract}

\tableofcontents

\section{Scope}

The purpose of this paper is not to introduce a new theorem or a new numerical phase contract. Its purpose is to freeze the comparison between:
\begin{itemize}
\item the live GitHub repository state;
\item the current local HAOS-IIP worktree;
\item the intended release scope for the next public synchronization push.
\end{itemize}

This is needed because the local repository has materially outrun the public GitHub branch. A clean synchronization release should say so explicitly rather than silently mixing comparison, numerics, foundations, and publish logistics into one undocumented push.

\section{Live GitHub Baseline}

At the time of comparison, the live remote branch is:
\begin{itemize}
\item repository: \texttt{tomislavrupic/HAOS-IIP};
\item branch: \texttt{main};
\item live commit: \texttt{cf5d48fb2178b92879a98dcdf0d52a4ec90ce60e};
\item live release headline: ``Freeze geometry\_emergence release 47.2''.
\end{itemize}

That means the public GitHub branch still reflects the geometry-emergence tranche and does not yet contain the later continuum and derivation stack.

\section{Local Delta Above the Live GitHub State}

Relative to live GitHub \texttt{47.2}, the local repository contains one coherent research-and-release tranche. It can be organized into four blocks.

\subsection{Continuum-scaling and operator updates}

The scalar branch is no longer limited to the older clean interior cubic check. The local repository now includes:
\begin{itemize}
\item perturbed point-set scalar controls;
\item reflected boundary-condition controls;
\item a strengthened scalar convergence note and updated result payload;
\item harmonic / coexact sector-aware transverse comparison;
\item full-window candidate and direct gap scans for the vector branch.
\end{itemize}

These artifacts collectively justify the bounded scalar statement that later became Paper \texttt{48.1}, while keeping the vector branch bounded by a real harmonic-floor obstruction rather than rhetorical optimism.

\subsection{HAOS-to-harmonic theorem ladder}

The largest conceptual change above live GitHub is the new foundations stack. The local repository now contains a bounded theorem ladder:
\[
T1 \to T2 \to T4 \to T5 \to T6 \to T7 \to T8,
\]
with the scalar toy bridge and locality cleanup notes filling important gaps.

This does not prove that HAOS explains all of harmonic mathematics. It does something narrower and operationally useful:
\begin{quote}
it shows how much Laplacian, Hodge, sector-decomposition, physical-sector, and refinement-stability structure can be derived from recoverable coherence under interaction before the argument becomes assumption rather than theorem.
\end{quote}

That bounded ladder is not present on live GitHub.

\subsection{Active-sector transport machinery}

The local repository also contains the first explicit \texttt{T8} execution tools. The derived physical branch is no longer compared only by separate spectra. Instead, restricted coexact modes are carried into a common probe space by explicit comparison maps \(J_n\), and then compared by:
\begin{itemize}
\item transported overlaps;
\item principal-angle cosines;
\item scaled-eigenvalue drift.
\end{itemize}

This is a genuine structural change in the repository because it moves the vector-sector refinement question from intuition to an explicit transported-branch identity test.

\subsection{Release papers already prepared locally}

Two synthesis papers exist locally but are not yet live on GitHub:
\begin{enumerate}
\item \texttt{48.1} Scalar Continuum Control, Harmonic-Coexact Sector Separation, and a Bounded Roadmap Toward Continuum Physics on HAOS-IIP;
\item \texttt{49.1} From Recoverable Coherence to Harmonic Operator Structure: A Bounded Derivation Program, Physical-Sector Restriction, and Continuum-Stability Roadmap for HAOS-IIP.
\end{enumerate}

This paper becomes the third release-layer artifact in the tranche:
\begin{enumerate}
\setcounter{enumi}{2}
\item \texttt{50.1} GitHub-State Comparison and Continuum-Foundations Synchronization Release for HAOS-IIP.
\end{enumerate}

\section{Detailed Comparison Map}

Table~\ref{tab:delta} compresses the comparison between the live GitHub state and the local repository state.

\begin{longtable}{p{0.23\linewidth}p{0.28\linewidth}p{0.38\linewidth}}
\caption{Repository delta from live GitHub \texttt{47.2} to the local synchronization tranche.}\\
\toprule
Area & Live GitHub state & Local state above GitHub \\
\midrule
\endfirsthead
\toprule
Area & Live GitHub state & Local state above GitHub \\
\midrule
\endhead
Scalar continuum & Geometry-emergence branch only & Strengthened scalar convergence controls with perturbed point sets, boundary conditions, updated note, and result payload \\
Vector-sector comparison & No bounded harmonic / coexact comparison tranche live & Harmonic/coexact split comparison, candidate-window scans, and direct harmonic-to-coexact gap scans \\
Foundations derivation & No public HAOS-to-harmonic theorem ladder & Bounded executed notes for \texttt{T1--T8} plus scalar-to-graded bridge and \texttt{F1} cleanup \\
Physical-sector transport & No public active-sector transport map & Explicit \(J_n\) transport runner, bounded transport note, mild-disorder extension through \(n=28\) \\
Paper spine & Ends at \texttt{47.2} & Adds prepared papers \texttt{48.1}, \texttt{49.1}, and the present \texttt{50.1} synchronization paper \\
Execution plan & No public bounded work-package note for the continuum tranche & Explicit execution-program note ordering \texttt{F1--F4}, \texttt{N1--N3}, \texttt{C1--C4} \\
\bottomrule
\label{tab:delta}
\end{longtable}

\section{What Is Stronger Now}

Three things are materially stronger than the public GitHub branch suggests.

\subsection{The scalar authority boundary}

The scalar branch now has a real bounded continuum-facing contract in the local kernel regime. This is not a grand continuum claim, but it is stronger than the older public state.

\subsection{The foundations authority boundary}

The repository no longer only \emph{uses} harmonic mathematics. It now contains a bounded argument for why harmonic operator structure is a natural normal form of recoverability on graded relational substrates under explicit HAOS-aligned assumptions.

\subsection{The vector diagnostic boundary}

The vector branch is still open, but it is open in a sharper way. The active coexact branch can now be tracked explicitly under refinement, and the clean baseline weakness can be separated from the mild-disorder and defect-supported cases rather than being left as an ambiguous nonconvergence statement.

\section{What Is Still Open}

This synchronization paper is only honest if it states the remaining open items clearly.

\subsection{Foundational open items}

The following are still not closed:
\begin{itemize}
\item the incidence-normalization caveat (\texttt{F2});
\item channel completeness and the no-extra-bridge rule (\texttt{F3});
\item the Hilbert-complex upgrade of the finite-dimensional sector theorem (\texttt{F4}).
\end{itemize}

\subsection{Vector-sector open items}

The vector branch still lacks:
\begin{itemize}
\item a frozen disorder-threshold result bundle;
\item full active-branch refinement closure in the clean baseline;
\item universality, effective-law, and geometry / response closure on the same transported branch.
\end{itemize}

The interrupted disorder-threshold work therefore remains active work, not release-level evidence.

\section{Release Scope for 50.1}

The intended release scope for \texttt{50.1} is:
\begin{itemize}
\item HAOS-IIP numerics associated with the continuum / derivation / transport tranche;
\item new foundations notes;
\item experiment notes, result bundles, and plots already cleanly stamped in the repository;
\item prepared papers \texttt{48.1} and \texttt{49.1};
\item the present synchronization paper and release-index updates.
\end{itemize}

The release should explicitly exclude:
\begin{itemize}
\item unrelated image assets;
\item any unfinished disorder-threshold outputs that were not cleanly frozen into the standard result path.
\end{itemize}

\section{Recommended Publish Action}

The correct next public action is not to pretend that GitHub is already current. It is to synchronize the repository honestly by publishing one coherent tranche that moves public state from \texttt{47.2} to the present bounded continuum / derivation / transport state.

That synchronized push should make the following true simultaneously:
\begin{enumerate}
\item GitHub contains the scalar continuum and vector-sector boundary work of Paper \texttt{48.1};
\item GitHub contains the bounded HAOS-to-harmonic derivation stack of Paper \texttt{49.1};
\item GitHub contains the active-sector transport tools and execution-program note;
\item GitHub contains the explicit repository-comparison record of Paper \texttt{50.1};
\item GitHub still does \emph{not} overclaim a finished vector-sector continuum closure.
\end{enumerate}

\section{Conclusion}

The live GitHub repository is still frozen at a geometry-emergence release. The local repository is already farther along. It contains a coherent continuum-facing tranche with stronger scalar evidence, a bounded HAOS-to-harmonic derivation stack, and the first explicit active-sector refinement transport machinery. That tranche is real and should be published. At the same time, the vector branch remains numerically unfinished, and the disorder-threshold question should remain outside the frozen claim set until a clean result bundle exists. Release \texttt{50.1} is therefore not a new physics claim. It is the synchronization document that makes the public state match the actual bounded local state.

\end{document}
